\documentclass[11pt]{extarticle}
\usepackage[margin=1.2in]{geometry}
\usepackage{amsmath,amssymb,amsthm}
\usepackage{../data_sharing_with_comp_ec26/style}
\newcommand{\timeind}{t}

\newcommand{\parahead}[1]{\vspace{0.1in}\noindent\textbf{#1. }}
\newcommand{\subparahead}[1]{\vspace{0.07in}\noindent\emph{#1. }}
\newcommand{\paraheadwnogap}[1]{\vspace{0.0in}\noindent\textbf{#1. }}
\newcommand{\subparaheadwnogap}[1]{\vspace{0.0in}\noindent\emph{#1. }}


\newcommand{\competitora}{a}
\newcommand{\competitorb}{b}
\newcommand{\arbitrarycompetitor}{z}

% Valuations and utilities
\newcommand{\numdataa}{n}
\newcommand{\numdatab}{m}

\newcommand{\numdatafora}{i}
\newcommand{\numdataforatt}[1]{\numdatafora_{#1}}
\newcommand{\numdataforat}{\numdataforatt{t}}

\newcommand{\numdataforb}{j}
\newcommand{\numdataforbtt}[1]{\numdataforb_{#1}}
\newcommand{\numdataforbt}{\numdataforbtt{t}}

\newcommand{\vala}{V^\competitora}
\newcommand{\valaii}[1]{\vala_{#1}}
\newcommand{\valai}{\valaii{i}}

\newcommand{\valb}{V^\competitorb}
\newcommand{\valbjj}[1]{\valb_{#1}}
\newcommand{\valbj}{\valbjj{j}}


\newcommand{\setexchanges}{\Scal}
\newcommand{\exchange}{s}
\newcommand{\exchangett}[1]{s_{#1}}
\newcommand{\exchanget}{\exchangett{t}}

\newcommand{\exchangea}{i}
\newcommand{\exchangeb}{j}

\newcommand{\utila}{U^\competitora}
\newcommand{\utilass}[1]{\utila_{#1}}
\newcommand{\utilas}{\utilass{s}}

\newcommand{\utilb}{U^\competitorb}
\newcommand{\utilbss}[1]{\utilb_{#1}}
\newcommand{\utilbs}{\utilbss{s}}

\newcommand{\prior}{D}
\newcommand{\priora}{\prior^a}
\newcommand{\priorb}{\prior^b}

\newcommand{\compcoeff}{\alpha}

\newcommand{\welfarefunc}{W}

\newcommand{\reward}{Y}
\newcommand{\rewardss}[1]{\reward_{#1}}
\newcommand{\rewards}{\rewardss{s}}

% Algorithm formulations
% \newcommand{\sharingalg}{\pi}
\newcommand{\algstatemap}{f}
\newcommand{\algstatemaptt}[1]{f_{#1}}
\newcommand{\algstatemapt}{\algstatemaptt{t}}

\newcommand{\algstopmap}{\delta}
\newcommand{\algstopmaptt}[1]{\algstopmap_{#1}}
\newcommand{\algstopmapt}{\algstopmaptt{t}}

\newcommand{\altalgstopmap}{\delta'}
\newcommand{\altalgstopmaptt}[1]{\altalgstopmap_{#1}}
\newcommand{\altalgstopmapt}{\altalgstopmaptt{t}}

\newcommand{\algendtime}{T} % Used to be \tau
\newcommand{\altalgendtime}{T'}

\newcommand{\algstate}{s} % Used to be \sigma
\newcommand{\algstatett}[1]{\algstate_{#1}}
\newcommand{\algstatet}{\algstatett{t}}

\newcommand{\timefilt}{\Gcal}
\newcommand{\timefilttt}[1]{\timefilt_{#1}}
\newcommand{\timefiltt}{\timefilttt{t}}

\newcommand{\statefilt}{\Fcal}
\newcommand{\statefiltss}[1]{\statefilt_{#1}}
\newcommand{\statefilts}{\statefiltss{s}}

\newcommand{\alg}{A}
\newcommand{\altalg}{A'}
\newcommand{\classalgs}{\Acal}
\newcommand{\classicalgs}{\Acal^{\text{IC}}}

\newcommand{\strategy}{\sigma}
\newcommand{\strategytt}[1]{\strategy_{#1}}
\newcommand{\strategyt}{\strategytt{t}}

\newcommand{\altstrategy}{\sigma'}
\newcommand{\altstrategytt}[1]{\altstrategy_{#1}}
\newcommand{\altstrategyt}{\altstrategytt{\timeind}}

\newcommand{\stoptime}{\tau}
\newcommand{\altstoptime}{\tau'}


\newcommand{\apolicy}{\pi} % for some reason \policy seems taken?
\newcommand{\setpolicies}{\Pi}

\newcommand{\exchangepolicy}{\pi}
\newcommand{\altexchangepolicy}{\pi'}

\newcommand{\setexchangepolicies}{\setpolicies^{\text{ex}}}
\newcommand{\seticexchangepolicies}{\setpolicies^{\text{IC}}}

% osas - stopping time paper
\newcommand{\osasreward}{Z}
\newcommand{\osasrewardss}[1]{\osasreward_{#1}}
\newcommand{\osasrewards}{\osasrewardss{s}}

\newcommand{\osasstrategy}{\sigma}
\newcommand{\osasstrategytt}[1]{\osasstrategy_{#1}}
\newcommand{\osasstrategyt}{\osasstrategytt{t}}

\newcommand{\osasstoptime}{\tau}

\newcommand{\osaspolicy}{\pi}
\newcommand{\osasstate}{s}
\newcommand{\osasstates}{S}

\newcommand{\osaspolicies}{\Pi}
\newcommand{\osaspoliciesss}[1]{\osaspolicies_{#1}}
\newcommand{\osaspoliciess}{\osaspoliciesss{s}}

\newcommand{\osassnellenv}{X}
\newcommand{\osassnellenvss}[1]{\osassnellenv_{#1}}
\newcommand{\osassnellenvs}{\osassnellenvss{s}}

\newcommand{\osasstatefilt}{\Fcal}
\newcommand{\osasstatefiltss}[1]{\osasstatefilt_{#1}}
\newcommand{\osasstatefilts}{\osasstatefiltss{s}}

\newcommand{\osasdirectsuccessorss}[1]{U(#1)}

\newcommand{\snellenv}{X}
\newcommand{\snellenvss}[1]{\snellenv_{#1}}
\newcommand{\snellenvs}{\snellenvss{s}}


\newcommand{\nextexchange}{u}

\newcommand{\directsuccessorss}[1]{U(#1)}

% \newcommand{\optalg}{\tilde{\alg}}
\newcommand{\optalg}{\alg}
% \newcommand{\optexpolicy}{\tilde{\apolicy}}
\newcommand{\optexpolicy}{\apolicy}

% \newcommand{\optstrategy}{\tilde{\sigma}}
\newcommand{\optstrategy}{\sigma}
\newcommand{\optstrategytt}[1]{\optstrategy_{#1}}
\newcommand{\optstrategyt}{\optstrategytt{t}}

% \newcommand{\optstoptime}{\tilde{\tau}}
\newcommand{\optstoptime}{\tau}

\newcommand{\zerotuple}{(0,0)}
% \newcommand{\zeroabrev}{}

\newcommand{\mbleexistsalgstatemap}{\tilde{\algstatemap}}
\newcommand{\mbleexistsalgstatemapt}{\tilde{\algstatemapt}}

% Notation for truthfulness proof
\newcommand{\opttruthexpolicy}{\pi'}
\newcommand{\opttruthexpolicykk}[1]{\opttruthexpolicy_{#1}}
\newcommand{\opttruthexpolicyk}{\opttruthexpolicykk{k}}

\newcommand{\opttruthstrategy}{\sigma'}
\newcommand{\opttruthstrategytt}[1]{\opttruthstrategy_{#1}}
\newcommand{\opttruthstrategyt}{\opttruthstrategytt{t}}

\newcommand{\opttruthstoptime}{\tau'}

\newcommand{\stopstateforpolicy}[1]{\alpha\rbr{#1}}

\newcommand{\fixedpath}{P}

\newcommand{\icsnellenv}{X}
\newcommand{\icsnellenvpp}[1]{\icsnellenv_{#1}}
\newcommand{\icsnellenvp}{\icsnellenvpp{\fixedpath}}

\newcommand{\rewardbound}{C}

% \newcommand{\policyfollowspath}[2]{#1\rbr{#2}}
\newcommand{\policyfollowspath}[2]{#2\rbr{#1}}
% \newcommand{\policyfollowspathpp}[2]{#2\rbr{#1}}

\newcommand{\exchangeindexrange}[2]{\exchangett{#1:#2}}
\newcommand{\strategyindexrange}[2]{\strategytt{#1:#2}}

\newcommand{\existspolicymatchingpath}[1]{E\rbr{#1}}

\newcommand{\validicextensions}{\Ucal}
\newcommand{\validicextensionspp}[1]{\validicextensions_{#1}}

\newcommand{\expolicymatchicopt}{\pi}
% \newcommand{\expolicymatchicopt}{\pi}

\newcommand{\lastfollowtime}{\xi}
\newcommand{\lastfollowtimeii}[1]{\lastfollowtime_{#1}}
\newcommand{\lastfollowtimei}{\lastfollowtimeii{i}}

\newcommand{\lastcommonfollowtime}{\chi}

\newcommand{\chosenpath}{\iota}
\usepackage[hidelinks]{hyperref}
\hypersetup{pdftitle={Existence of an IC Optimal Resource Exchange Policy}}
\declaretheorem[name=Theorem]{theorem}
\declaretheorem[name=Lemma,sibling=theorem]{lemma}
\declaretheorem[name=Proposition,sibling=theorem]{proposition}
\declaretheorem[name=Remark,sibling=theorem]{remark}

\title{Existence of an IC Optimal Resource Exchange Policy}
\author{}
\date{November 13, 2026}

\begin{document}
\maketitle
\vspace{-1.5em}

\begin{abstract}
An incentive-compatible optimal policy always exists in the resource exchange
model with finitely many goods and bounded measurable utilities and welfare.
The policy is deterministic, and valuations may have arbitrary distributions,
including atoms and continuous components. The proof keeps track of both
agents' expected continuation utilities together with expected welfare.
Their attainable triples form a compact set. A measurable selection theorem
allows every triple used in the construction to be implemented exactly.
No closure-under-pasting assumption or countability of the policy class is
required. Incentive compatibility here concerns early exit, with truthful
observation of valuations, as in the original model.
\end{abstract}

\section{Model and result}

Let $n,m\in\NN_{\geq 0}$ be finite. The random valuation vectors
$(\valaii{1},\ldots,\valaii{n})$ and
$(\valbjj{1},\ldots,\valbjj{m})$ take values in $\RR^n$ and $\RR^m$,
respectively, and have given independent distributions. Coordinates within
each vector may be arbitrarily dependent. Set $\valaii{0}=\valbjj{0}=0$.
An exchange is an element
\[
    \setexchanges=\cbr{0,\ldots,n}\times\cbr{0,\ldots,m}.
\]
The order on $\setexchanges$ is coordinatewise: $u>s$ means that both
coordinates weakly increase and at least one strictly increases.
At $s=(i,j)$, the publicly observed information is the full prefix
\[
    H_s=\rbr{\valaii{1},\ldots,\valaii{i},
                    \valbjj{1},\ldots,\valbjj{j}},
    \qquad \statefilts=\sigma(H_s).
\]
The prefix space is $\Hcal_s=\RR^i\times\RR^j$, with $\Hcal_0$ a
singleton and $0=(0,0)$. In particular, $\Fcal_0$ is trivial.

Let $\utila,\utilb,\welfarefunc:\RR^2\to\RR$ be bounded Borel
functions, with $\utila(0,0)=\utilb(0,0)=\welfarefunc(0,0)=0$.
Terminal utilities and welfare at $s=(i,j)$ are
\[
    \utilas=\utila(\valai,\valbj),\qquad
    \utilbs=\utilb(\valai,\valbj),\qquad
    \rewards=\welfarefunc(\utilas,\utilbs).
\]
All measurability below is Borel measurability on these finite-dimensional
observation spaces; conditional expectation identities are almost sure.

A deterministic exchange policy $\pi=(\sigma,\tau)$ starts at
$\sigma_0=0$. At each round it uses the observations so far to stop or
choose a next exchange $\sigma_{t+1}>\sigma_t$. It must stop at $(n,m)$.
Let $\Gcal_t$ be the information generated by the observations and exchange
history through round $t$. Decisions are measurable with respect to this
information. Thus $\tau\leq N:=n+m$. After termination we freeze the state
and information; this convention does not affect any payoff or early-exit
decision. Section~\ref{sec:information} verifies that history-based rules
used below are representable by the original model's measurable maps of the
observed prefixes.

A policy is \emph{incentive compatible} (IC) if, for each $z\in\{a,b\}$
and every $\Gcal_t$-stopping time $\rho$,
\begin{equation}
    \EE\sbr{U^z_{\sigma_{\tau\wedge\rho}}}
    \leq \EE\sbr{U^z_{\sigma_\tau}}.
    \label{eq:IC}
\end{equation}
It suffices to consider $0\leq\rho\leq N$.
Write $\seticexchangepolicies$ for the class of such policies.
An \emph{IC optimal} policy maximizes $\EE[Y_{\sigma_\tau}]$ over this
class. This is the optimum subject to IC, rather than a claim that an
unconstrained welfare optimum is IC.

\begin{theorem}[Existence and compactness]\label{thm:main}
In every environment just defined, the set
\begin{equation}
    \Kcal^{\mathrm{IC}}
    =\cbr{
       \rbr{\EE[\utilass{\sigma_\tau}],
             \EE[\utilbss{\sigma_\tau}],
             \EE[\rewardss{\sigma_\tau}]}:
       (\sigma,\tau)\in\seticexchangepolicies
      }
    \subseteq\RR^3
    \label{eq:attainable}
\end{equation}
is nonempty and compact. Consequently, there is a deterministic measurable
policy $\pi^*=(\sigma^*,\tau^*)\in\seticexchangepolicies$ such that
\[
    \EE\sbr{Y_{\sigma^*_{\tau^*}}}
    =\max_{\pi=(\sigma,\tau)\in\seticexchangepolicies}
      \EE\sbr{Y_{\sigma_\tau}}.
\]
This policy is also individually rational:
$\EE[U^a_{\sigma^*_{\tau^*}}]\geq0$ and
$\EE[U^b_{\sigma^*_{\tau^*}}]\geq0$.
No further assumptions on the valuation distributions, utility functions,
or welfare function are needed.
\end{theorem}

The proof below is independent of the original paper's results. Its
measure-theoretic input is a standard parameterized selection theorem,
stated explicitly in Section~\ref{sec:selection} with a precise reference.
All exchange-policy and incentive arguments are proved here.

\section{Information and the conditional form of IC}
\label{sec:information}

For a deterministic policy, a prescribed finite exchange path ending at
$s$ is determined by $H_s$: all earlier observed prefixes are coordinate
projections of $H_s$, and the policy's decisions can be simulated from
those prefixes. There are only finitely many strictly increasing paths.
It follows that the events
\begin{align*}
    R_s&=\cbr{\text{the policy visits }s\text{ at or before termination}},\\
    B_{s,u}&=\cbr{\text{the policy visits }s\text{ and next moves to }u}
\end{align*}
belong to $\Fcal_s$ for every $u>s$. The event of stopping at $s$ also
belongs to $\Fcal_s$. Likewise, on $\{t\leq\tau,\sigma_t=s\}$ the trace
of $\Gcal_t$ is the trace of $\Fcal_s$: the current prefix contains all
earlier observations, and the earlier actions are reconstructible from it.

This argument also applies to a rule with deterministic internal memory
updated measurably from observations. Starting from its fixed initial
memory, simulate the rule to recover the memory and past exchanges at
round $t$ from the current prefix. On an input inconsistent with this
simulation, define the rule to stop. Simulation takes at most $N$ steps
and uses only Borel operations and a finite collection of states. The
resulting stop/next-exchange maps are therefore Borel functions of the
round and observed prefix, as required in the original algorithm
definition. In particular, such memory neither supplies a random seed nor
reveals unobserved valuations.

For each $s<u$, fix a Borel regular conditional distribution
\begin{equation}
    Q_{s,u}(h,dh')=P(H_u\in dh'\mid H_s=h).
    \label{eq:kernel}
\end{equation}
These kernels exist because the prefix spaces are standard Borel spaces.
At $s=0$, the kernel is the unconditional law of $H_u$.
Conditioning on reaching $s$, or on the choice of $u$ at $s$, does not
change the conditional law in~\eqref{eq:kernel}: those events are already
$\Fcal_s$-measurable. The same holds after conditioning on any
deterministic internal memory reconstructible from $H_s$.

We use the kernels only in almost-sure identities along actual policy
paths. Their arbitrary values on exceptional histories need not satisfy
any pointwise consistency identities. Indeed, if $D\subseteq\Hcal_s$
has $P(H_s\in D)=0$, then
$P(R_s\cap\{H_s\in D\})=0$. There are finitely many states and transitions,
so exceptional sets cannot become relevant through adaptive choices.

\begin{lemma}[Conditional participation constraints]\label{lem:local}
A policy $\pi=(\sigma,\tau)$ is IC if and only if, for every
$t\in\{0,\ldots,N\}$ and $z\in\{a,b\}$,
\begin{equation}
    \EE\sbr{U^z_{\sigma_\tau}\mid\Gcal_t}
    \geq U^z_{\sigma_t}
    \quad\text{almost surely on }\{t\leq\tau\}.
    \label{eq:local}
\end{equation}
\end{lemma}
\begin{proof}
At $t=\tau$ there is equality. If~\eqref{eq:local} fails on a set of
positive probability, then for some $t,z$ the event
\[
    D=\{t<\tau\}\cap
       \cbr{U^z_{\sigma_t}>
               \EE[U^z_{\sigma_\tau}\mid\Gcal_t]}
       \in\Gcal_t
\]
has positive probability. Set $\rho=t$ on $D$ and $\rho=N$ otherwise.
This is a stopping time, and the expected gain from this deviation is
\[
    \EE\sbr{1_D\rbr{U^z_{\sigma_t}-U^z_{\sigma_\tau}}}
    =\EE\sbr{1_D\rbr{U^z_{\sigma_t}
             -\EE[U^z_{\sigma_\tau}\mid\Gcal_t]}}>0,
\]
contradicting IC.

Conversely, if~\eqref{eq:local} holds, any stopping time $\rho\leq N$
satisfies
\begin{align*}
    \EE\sbr{U^z_{\sigma_\tau}
                 -U^z_{\sigma_{\tau\wedge\rho}}}
    &=\sum_{t=0}^{N}
      \EE\sbr{1_{\{\rho=t<\tau\}}
                  \rbr{U^z_{\sigma_\tau}-U^z_{\sigma_t}}}\\
    &=\sum_{t=0}^{N}
      \EE\sbr{1_{\{\rho=t<\tau\}}
                  \rbr{\EE[U^z_{\sigma_\tau}\mid\Gcal_t]
                                  -U^z_{\sigma_t}}}\geq0.
\end{align*}
This is~\eqref{eq:IC} for both agents.
\end{proof}

\section{The measurable-selection input}\label{sec:selection}

We use the following identity-map case of the measurable
``measurable choice'' theorem; see
Jacobovic, Levy, and Solan~\cite[Theorem 4, pp.~226--227]{JLS}.
Here a correspondence has a Borel graph if its set of input--output
pairs is Borel. No atomlessness assumption is made.

\begin{theorem}[Parameterized selection integral, quoted]\label{thm:selection}
Let $X,Z$ be standard Borel spaces, $C\subseteq\RR^d$ compact,
$F:X\times Z\rightrightarrows C$ have a Borel graph and nonempty compact
values, and $Q$ be a Borel probability kernel from $X$ to $Z$. Define
\[
    I(x)=\cbr{\int f(z)\,Q(x,dz):
                 f:Z\to C\text{ Borel},\ f(z)\in F(x,z)\text{ for all }z}.
\]
Then $I$ has a Borel graph and nonempty compact values. There is a Borel
map $g:\operatorname{Gr}(I)\times Z\to C$ satisfying, for every
$(x,v)\in\operatorname{Gr}(I)$,
\[
    g(x,v,z)\in F(x,z)\quad\text{for all }z,\qquad
    \int g(x,v,z)\,Q(x,dz)=v.
\]
\end{theorem}

The integral here is the set of integrals of actual measurable
selections, not its closure or convex hull. Thus its compactness concerns
exact attainable vectors. The last assertion selects a realizing
function jointly measurably in both the observed history and the desired
vector. This is the step needed to implement a continuum of continuation
choices by one admissible policy. Applying the theorem with $d=3$ will
preserve both promised utilities and welfare exactly.

\section{Backward construction of compact payoff sets}

Choose $M<\infty$ bounding the absolute values of $\utila,\utilb$, and
$\welfarefunc$, and set $C=[-M,M]^3$. At $h\in\Hcal_s$, write
\[
    r_s(h)=\rbr{u_s^a(h),u_s^b(h),y_s(h)}\in C
\]
for the three payoffs from stopping at $s$. These are the Borel functions
of the prefix specified by the model; in particular,
$r_s(H_s)=(U_s^a,U_s^b,Y_s)$ and $r_0=(0,0,0)$.
For a triple $v\in\RR^3$, label its coordinates $v^a,v^b,v^w$.

Define correspondences $K_s:\Hcal_s\rightrightarrows C$ in reverse
order of $i+j$, for $s=(i,j)$. At the maximal state, let
\[
    K_{(n,m)}(h)=\{r_{(n,m)}(h)\}.
\]
Once $K_u$ is defined for every $u>s$, set
\begin{align}
    I_{s,u}(h)
    &=\cbr{\int f(h')\,Q_{s,u}(h,dh'):
        f:\Hcal_u\to C\text{ Borel},\ f(h')\in K_u(h')\text{ for all }h'},
       \label{eq:I}\\
    D_{s,u}(h)
    &= I_{s,u}(h)\cap
       \cbr{v\in C:v^a\geq u_s^a(h),\ v^b\geq u_s^b(h)},
       \label{eq:D}\\
    K_s(h)
    &=\{r_s(h)\}\ \cup\!\bigcup_{u>s}D_{s,u}(h).
       \label{eq:K}
\end{align}
The union is over \emph{all} strict successors, including exchanges that
give several additional goods at once.

\begin{lemma}\label{lem:compact}
For every $s$, $K_s$ has a Borel graph and nonempty compact values.
Every $v\in K_s(h)$ satisfies
\begin{equation}
    v^a\geq u_s^a(h),\qquad v^b\geq u_s^b(h).
    \label{eq:coordinate-bounds}
\end{equation}
Moreover, for each $s<u$ there is a Borel map
$g_{s,u}:\operatorname{Gr}(I_{s,u})\times\Hcal_u\to C$ such that
\begin{equation}
    g_{s,u}(h,v,h')\in K_u(h'),\qquad
    \int g_{s,u}(h,v,h')\,Q_{s,u}(h,dh')=v.
    \label{eq:realizer}
\end{equation}
\end{lemma}
\begin{proof}
The assertions about $K_s$ are immediate at the maximal state. Inductively,
apply Theorem~\ref{thm:selection} with
\[
    X=\Hcal_s,\quad Z=\Hcal_u,\quad d=3,\quad
    F(h,h')=K_u(h'),\quad Q=Q_{s,u}.
\]
The inductive hypotheses verify all its assumptions. Thus $I_{s,u}$ has
a Borel graph and compact values, and~\eqref{eq:realizer} holds.
Intersecting each section with the two closed halfspaces
in~\eqref{eq:D} preserves compactness (possibly giving the empty set).
Its graph is Borel because the two thresholds are Borel functions of
$h$. The graph of $h\mapsto\{r_s(h)\}$ is Borel. The finite union
in~\eqref{eq:K} therefore has a Borel graph and compact values, and is
nonempty because it contains $r_s(h)$.
The coordinate bounds follow directly from~\eqref{eq:D} and equality
for the stopping triple. This completes the induction.
\end{proof}

\section{Every IC policy has its payoff triple in the initial set}

\begin{proposition}\label{prop:upper}
$\Kcal^{\mathrm{IC}}\subseteq K_0$.
\end{proposition}
\begin{proof}
Fix an IC policy $\pi$ and let its terminal payoff vector be
\[
    Z=\rbr{U^a_{\sigma_\tau},U^b_{\sigma_\tau},Y_{\sigma_\tau}}.
\]
For every $s$, choose a Borel version $v_s:\Hcal_s\to C$ of its
conditional mean:
\begin{equation}
    v_s(H_s)=\EE[Z\mid\Fcal_s].
    \label{eq:v}
\end{equation}
These versions exist by conditional expectation and measurable
factorization through $H_s$; values outside $C$ on a null set can be
replaced by $r_s$. We prove, by reverse induction on $i+j$, that
\begin{equation}
    v_s(H_s)\in K_s(H_s)\quad\text{almost surely on }R_s.
    \label{eq:membership}
\end{equation}

At a history where the policy stops at $s$, its terminal vector is
$r_s(H_s)$. Since this stopping event is in $\Fcal_s$, conditional
expectation gives $v_s(H_s)=r_s(H_s)$ there. This proves the base case
$s=(n,m)$ and handles every stopping event in the induction.

Suppose~\eqref{eq:membership} is known at all $u>s$. Choose a Borel
set $A_u\subseteq\Hcal_u$ representing $R_u=\{H_u\in A_u\}$ and put
\[
    f_u(h')=
    \begin{cases}
       v_u(h'),&h'\in A_u\text{ and }v_u(h')\in K_u(h'),\\
       r_u(h'),&\text{otherwise}.
    \end{cases}
\]
This is Borel by Lemma~\ref{lem:compact}, belongs to $K_u(h')$ at
\emph{every} $h'$, and satisfies
$f_u(H_u)=v_u(H_u)$ almost surely on $R_u$ by the induction hypothesis.

On the event $B_{s,u}$, the policy necessarily reaches $u$. Since
$B_{s,u}\in\Fcal_s$ and $B_{s,u}\subseteq R_u$, conditional expectation
and the tower property yield, on $B_{s,u}$ almost surely,
\begin{align}
    \int f_u(h')\,Q_{s,u}(H_s,dh')
    &=\EE[f_u(H_u)\mid\Fcal_s]\notag\\
    &=\EE[v_u(H_u)\mid\Fcal_s]
      =\EE[Z\mid\Fcal_s]=v_s(H_s).
    \label{eq:tower}
\end{align}
Here $\Fcal_s\subseteq\Fcal_u$ because $s<u$. The integral on the
left belongs to $I_{s,u}(H_s)$ by~\eqref{eq:I}.

Lemma~\ref{lem:local} and the information identity in
Section~\ref{sec:information} also imply
\[
    v_s^a(H_s)\geq U_s^a,\qquad
    v_s^b(H_s)\geq U_s^b
    \quad\text{almost surely on }R_s.
\]
To see that this applies at a state reached at a random round, partition
$R_s$ into its finitely many events $\{t\leq\tau,\sigma_t=s\}$ and use
the equality of the trace information fields on each event.
Consequently,~\eqref{eq:tower} places $v_s(H_s)$ in $D_{s,u}(H_s)$
on $B_{s,u}$. The stopping event and the finitely many $B_{s,u}$
partition $R_s$, completing the induction.

Finally, $R_0$ is the whole probability space and $\Fcal_0$ is trivial.
Thus~\eqref{eq:membership} at $0$ says $\EE[Z]\in K_0$.
\end{proof}

\section{Exact deterministic implementation of the initial set}

\begin{proposition}\label{prop:lower}
Every $v\in K_0$ is the expected payoff triple of a deterministic
measurable IC exchange policy. Hence $K_0\subseteq\Kcal^{\mathrm{IC}}$.
\end{proposition}
\begin{proof}
Fix $v\in K_0$. We construct a policy carrying an internal target vector
$q_t$, initialized by $q_0=v$. Fix an ordering of the finite set of
states. At an active history with current state $s$, observed prefix
$h$, and target $q\in K_s(h)$, proceed as follows.
\begin{enumerate}
    \item If $q=r_s(h)$, stop.
    \item Otherwise choose the first $u>s$ in the fixed ordering for
    which $q\in D_{s,u}(h)$. Such a $u$ exists by~\eqref{eq:K}.
    Make this exchange and, after observing $h'=H_u$, set
    \begin{equation}
        q_{t+1}=g_{s,u}(h,q,h').
        \label{eq:update}
    \end{equation}
\end{enumerate}
All choices are Borel: equality with $r_s(h)$ is a Borel condition,
and membership in each $D_{s,u}(h)$ is Borel. Choosing the first
eligible member of a finite list preserves measurability. The update
is Borel by~\eqref{eq:realizer}. Its target lies in $K_u(h')$ at every
step. At $(n,m)$ the target must equal the stopping triple, so the
policy stops after at most $N$ moves. The deterministic memory is
representable by the original prefix-based policy maps, by the
simulation argument of Section~\ref{sec:information}.

After stopping, keep $q_t=q_\tau$. We next prove that this bounded
vector process is a martingale for the actual observation filtration.
At a continuing round, condition on the current history. The selected
$s,u,h,q$ are then known, and the conditional law of the next observed
prefix is $Q_{s,u}(h,\cdot)$. Therefore~\eqref{eq:realizer} gives
\begin{equation}
    \EE[q_{t+1}\mid\Gcal_t]
    =\int g_{s,u}(h,q,h')\,Q_{s,u}(h,dh')
    =q_t.
    \label{eq:martingale}
\end{equation}
The identity also holds after stopping because the process is constant.
The kernel identity with the random parameters $(h,q)$ is valid since
the integrand is jointly Borel and these parameters are determined by
the current history. As explained in Section~\ref{sec:information},
the kernels need only be valid almost surely at reached histories.

At termination the target equals the realized terminal payoff vector:
\[
    q_\tau=r_{\sigma_\tau}(H_{\sigma_\tau})
       =\rbr{U^a_{\sigma_\tau},U^b_{\sigma_\tau},Y_{\sigma_\tau}}.
\]
Iterating~\eqref{eq:martingale} to the deterministic horizon $N$ gives
\begin{equation}
    q_t=\EE[q_\tau\mid\Gcal_t],\qquad
    \EE[q_\tau]=q_0=v.
    \label{eq:promise}
\end{equation}
Thus the target is precisely the conditional expected terminal payoff
vector, and the initial target is attained exactly.

Finally, $q_t\in K_{\sigma_t}(H_{\sigma_t})$ at every reached state,
so~\eqref{eq:coordinate-bounds} and~\eqref{eq:promise} imply
\[
    \EE[U^z_{\sigma_\tau}\mid\Gcal_t]
       =q_t^z\geq U^z_{\sigma_t}
       \quad\text{on }\{t\leq\tau\},\qquad z\in\{a,b\}.
\]
Lemma~\ref{lem:local} proves IC against every alternative stopping time.
\end{proof}

\begin{proof}[Proof of Theorem~\ref{thm:main}]
Propositions~\ref{prop:upper} and~\ref{prop:lower} give
$\Kcal^{\mathrm{IC}}=K_0$. Lemma~\ref{lem:compact} makes this set
nonempty and compact. The continuous coordinate function $v\mapsto v^w$
therefore attains a maximum at some $v^*\in K_0$.
Proposition~\ref{prop:lower} implements $v^*$ by a deterministic IC
policy. Proposition~\ref{prop:upper} bounds the welfare of every IC
policy by $(v^*)^w$, proving optimality. At the initial state,
\eqref{eq:coordinate-bounds} gives $(v^*)^a,(v^*)^b\geq0$, proving IR.
Equivalently, IR follows directly by choosing $\rho=0$ in the IC
definition. If $n=m=0$, the construction consists solely of stopping
at $0$, and the same conclusion holds.
\end{proof}

\section{Scope and interpretation}

The recursion retains all attainable continuation triples because a
welfare-maximizing continuation at a later history can lower an agent's
expected utility promised earlier. Equations~\eqref{eq:realizer}
and~\eqref{eq:promise} preserve both utility coordinates, so the
construction maintains earlier participation incentives while attaining
the selected welfare. It does not require arbitrary pasting of IC
policies to preserve IC.

The conclusion covers arbitrary continuous, discrete, and mixed
valuation distributions and bounded discontinuous payoffs. Atoms are
allowed in Theorem~\ref{thm:selection}; the sets in~\eqref{eq:I} need
not be convex. No public or private randomization is added to the
policy. The proof also does not use independence between the agents'
valuation vectors: it remains valid for any joint law on
$\RR^n\times\RR^m$ with the same public-prefix information structure.

This is an existence result and an exact set-valued characterization.
Computing the measurable selection integrals for general continuous
priors is a separate algorithmic problem. The result addresses the
early-exit IC optimum defined in the paper; strategic misreporting would
require additional incentive constraints and is not part of this model.

\begin{thebibliography}{9}
\bibitem{JLS}
Royi Jacobovic, Yehuda John Levy, and Eilon Solan.
\newblock Bayesian games with nested information.
\newblock \emph{Theoretical Economics}, 21 (2026), 204--240.
\newblock Theorem 4, pp.~226--227 (identity-map case).
\newblock \href{https://doi.org/10.3982/TE6461}{doi:10.3982/TE6461}.
\newblock \href{https://eprints.gla.ac.uk/362558/2/362558.pdf}{Open-access text}.
\end{thebibliography}

\end{document}
